\documentclass{amsart}
\usepackage{amsmath,amssymb,amsthm,hyperref,bm}
\newtheorem{theorem}{Theorem}
\newtheorem{proposition}{Proposition}
\newtheorem{lemma}{Lemma}
\newtheorem{definition}{Definition}
\newtheorem{corollary}{Corollary}
\theoremstyle{remark}
\newtheorem{remark}{Remark}

\begin{document}

\title{Construct validity of team psychological safety:
a measurement-theoretic treatment}
\maketitle

\begin{abstract}
We give a complete treatment of the three requested tasks. Task~1 is a
\emph{conceptual} specification, which we make precise by fixing a latent-variable
data-generating model and locating each construct as a distinct latent dimension
together with its level of analysis, referent, and time horizon. Tasks~2 and~3 are
\emph{empirical/statistical}: their content is a set of decision rules on observable
statistics. We state each rule as a precise proposition and prove the
mathematical facts on which the rules rest---identifiability and nesting of the
confirmatory-factor (CFA) models, validity of the $\chi^2$-difference test for
discriminant validity, the Spearman--Brown relation linking $\mathrm{ICC}(1)$ to
$\mathrm{ICC}(2)$, the meaning of the within-group agreement index $r_{WG}$, and the
incremental-validity (hierarchical regression) criterion. We then state exactly
which inputs are empirical (i.e.\ must be read off a data set, as in Edmondson's
1999 field study) and which conclusions follow deductively once those inputs are
in hand. The role of the empirical inputs is made explicit so that the logical
status of every claim is unambiguous.
\end{abstract}

\section{Scope and logical status of the problem}\label{sec:scope}

The problem is cross-disciplinary: it asks us to \emph{establish} the validity of a
psychological construct. ``Construct validity'' is not a theorem with a proof from
axioms; it is a verdict reached by checking that a battery of statistical criteria
hold on data. Accordingly we split the problem into two logically distinct parts.

\begin{itemize}
\item[(C)] \textbf{Conceptual claims} (Task~1). These are definitions and the
relations between them. They are settled by exhibiting a model in which the
constructs are formally distinct objects; this is what Section~\ref{sec:concept}
does.
\item[(E)] \textbf{Empirical criteria} (Tasks~2--3). Each criterion is a predicate
on observable statistics (reliabilities, agreement indices, fit indices,
regression coefficients). We prove the \emph{mathematical} content of each
criterion: that the statistic is well defined, that it has the stated sampling
behaviour, and that, \emph{conditional on the data exhibiting a stated pattern},
the criterion delivers the stated conclusion. The numerical values that decide
whether the pattern holds are observations, not theorems; we flag every such value
as an \emph{empirical input} (EI).
\end{itemize}

This separation is the honest one: no amount of mathematics can prove from first
principles that, in the population of real teams, psychological safety happens to
have a factor correlation of (say) $0.6$ rather than $1.0$ with cohesiveness. What
mathematics \emph{can} do, and what we do below, is prove that the test procedures
used to detect such a separation are valid, and that the conjunction of the
criteria entails the discriminant-validity conclusion. The empirical instantiation
we reference throughout is Edmondson (1999, \emph{Administrative Science
Quarterly} 44:350--383), whose field data supply the EI values; our contribution
is the rigorous skeleton.

\section{Task 1: Conceptual specification and distinctness}\label{sec:concept}

\subsection{A common formal model}
Fix a population of teams indexed by $j=1,\dots,J$, team $j$ having $n_j\ge 2$
members indexed $i=1,\dots,n_j$. Let
\[
\Omega=\{\,\text{PS},\ \text{COH},\ \text{TR},\ \text{SAT},\ \text{EFF}\,\}
\]
denote the five candidate constructs: psychological safety (PS), group
cohesiveness (COH), interpersonal trust (TR), work-group satisfaction (SAT), and
team efficacy (EFF). Each construct $c\in\Omega$ is modelled as a latent random
variable carried at a specified \emph{level of analysis} $L(c)$, with a specified
\emph{referent} $R(c)$ and \emph{time horizon} $H(c)$. The data-generating model is
a multilevel reflective measurement model: for member $i$ of team $j$ and
indicator $m$ of construct $c$,
\begin{equation}\label{eq:cfa}
x^{(c)}_{ijm}=\nu^{(c)}_m+\lambda^{(c)}_m\,\eta^{(c)}_{ij}+\varepsilon^{(c)}_{ijm},
\qquad
\eta^{(c)}_{ij}=\theta^{(c)}_{j}+u^{(c)}_{ij},
\end{equation}
where $\nu^{(c)}_m$ is an intercept, $\lambda^{(c)}_m\neq0$ a loading,
$\theta^{(c)}_{j}$ a team-level latent score, $u^{(c)}_{ij}$ a within-team
deviation, and $\varepsilon^{(c)}_{ijm}$ measurement error, all mutually
independent with $\mathbb{E}[u]=\mathbb{E}[\varepsilon]=0$ and finite variances
$\sigma^2_{u,c}$, $\sigma^2_{\varepsilon,c}$. Write $\theta_j=(\theta^{(c)}_j)_{c\in\Omega}$
for the vector of team-level latent scores, with population correlation matrix
$\Phi=(\phi_{cc'})_{c,c'\in\Omega}$, $\phi_{cc}=1$.

This single model lets us state distinctness precisely.

\begin{definition}[Conceptual distinctness]\label{def:distinct}
Constructs $c$ and $c'$ are \emph{conceptually distinct} if they differ in at least
one of the three coordinates $(L,R,H)$, and \emph{empirically distinct} if in
addition $|\phi_{cc'}|<1$ (Section~\ref{sec:disc} makes the test of this precise).
\end{definition}

\subsection{The coordinates of psychological safety}
We fix:
\[
L(\mathrm{PS})=\text{team (shared belief)},\quad
R(\mathrm{PS})=\text{interpersonal risk-taking climate},\quad
H(\mathrm{PS})=\text{enduring/quasi-stable}.
\]
\emph{Level.} PS is a team-level shared belief: it is the construct
$\theta^{(\mathrm{PS})}_j$, well defined only when within-team variation
$\sigma^2_{u,\mathrm{PS}}$ is small relative to between-team variation
$\sigma^2_{\theta,\mathrm{PS}}:=\operatorname{Var}(\theta^{(\mathrm{PS})}_j)$
(operationalised by $r_{WG}$, $\mathrm{ICC}(1)$, $\mathrm{ICC}(2)$ in
Section~\ref{sec:agg}). This contrasts with TR, whose primitive referent is
\emph{dyadic} ($i\to i'$): trust is a property of an ordered pair, and a team-level
trust score is a (lossy) aggregate of a dyadic array, not a primitive shared
belief.

\emph{Referent.} The referent of PS is the perceived \emph{consequences of
interpersonal risk} (will I be punished for speaking up, admitting error, asking a
question?). This differs from:
\begin{itemize}
\item COH, whose referent is interpersonal \emph{attraction}/desire to remain in
the group --- a valenced bond, not a risk appraisal; high cohesion can even
\emph{lower} the willingness to voice dissent (groupthink), so COH and PS can move
oppositely, which is incompatible with $\phi_{\mathrm{PS},\mathrm{COH}}=1$.
\item TR, whose referent is \emph{reliance} on a specific other's benevolence or
competence (a dyadic expectation about another agent), whereas PS is a
\emph{collective} appraisal of the team \emph{climate} and need not be reducible to
the average of dyadic trusts.
\item SAT, whose referent is \emph{contentment} with the group as an object of
affect --- an evaluative attitude, not a belief about the safety of action.
\item EFF, whose referent is the team's \emph{capability} to execute tasks
(``can we perform?''), an outcome-expectancy about task performance, orthogonal in
referent to the social-risk appraisal of PS (``is it safe to take interpersonal
risk?'').
\end{itemize}

\emph{Time horizon.} PS is an enduring, quasi-stable climate perception (it should
exhibit test--retest stability, Section~\ref{sec:rel}), distinguishing it from
transient state affect; SAT, by contrast, is more state-like and reactive to recent
events. Thus on coordinate $H$, $\mathrm{PS}\neq\mathrm{SAT}$ as well.

\begin{proposition}[Conceptual distinctness, Task~1]\label{prop:concept}
Under the coordinate assignments above, PS is conceptually distinct
(Definition~\ref{def:distinct}) from each of COH, TR, SAT, EFF.
\end{proposition}
\begin{proof}
We exhibit a differing coordinate for each pair. $(\mathrm{PS},\mathrm{TR})$:
differ in $L$ (team shared belief vs.\ dyadic) and $R$ (risk climate vs.\
reliance). $(\mathrm{PS},\mathrm{COH})$: differ in $R$ (risk climate vs.\
attraction). $(\mathrm{PS},\mathrm{SAT})$: differ in $R$ (belief about safety of
action vs.\ contentment) and $H$ (enduring climate vs.\ state-reactive).
$(\mathrm{PS},\mathrm{EFF})$: differ in $R$ (interpersonal-risk appraisal vs.\
task-capability expectancy). In every case at least one of $(L,R,H)$ differs, which
is the definition. \end{proof}

\section{Task 2: A measurement instrument and its validation}\label{sec:meas}

\subsection{The instrument}
The instrument is a $p$-item reflective scale $(x_{ijm})_{m=1}^{p}$ for PS
following \eqref{eq:cfa} (e.g.\ Edmondson's 7-item scale, with appropriate items
reverse-scored so all $\lambda^{(\mathrm{PS})}_m>0$). The member-level scale score
is $s_{ij}=\frac1p\sum_{m=1}^{p}x^{(\mathrm{PS})}_{ijm}$ and the team-level score is
the within-team mean $\bar s_{j}=\frac1{n_j}\sum_{i=1}^{n_j}s_{ij}$. Parallel
scales are administered for COH, TR, SAT, EFF and for the outcomes (team learning
behaviour, error detection, supervisor-rated performance).

\subsection{(2a) Reliability}\label{sec:rel}
\begin{proposition}[Internal consistency lower bound]\label{prop:alpha}
Under \eqref{eq:cfa} with equal loadings $\lambda^{(\mathrm{PS})}_m\equiv\lambda$
and uncorrelated errors (tau-equivalence), Cronbach's $\alpha$ equals the
composite reliability
$\omega=\dfrac{(\sum_m\lambda_m)^2\operatorname{Var}(\eta)}
{(\sum_m\lambda_m)^2\operatorname{Var}(\eta)+\sum_m\operatorname{Var}(\varepsilon_m)}$;
in general $\alpha\le\omega\le 1$, so $\alpha$ is a lower bound for the proportion
of true-score variance.
\end{proposition}
\begin{proof}
Let $X=\sum_{m=1}^{p}x_m$ with $x_m=\lambda_m\eta+\varepsilon_m$ (suppressing
indices). Then $\operatorname{Var}(X)=(\sum_m\lambda_m)^2\sigma^2_\eta
+\sum_m\sigma^2_{\varepsilon_m}$ and the true-score (common) variance is
$\sigma^2_T=(\sum_m\lambda_m)^2\sigma^2_\eta$, giving $\omega=\sigma_T^2/\operatorname{Var}(X)$.
Cronbach's $\alpha=\frac{p}{p-1}\big(1-\frac{\sum_m\operatorname{Var}(x_m)}{\operatorname{Var}(X)}\big)$.
A classical inequality (Novick--Lewis 1967) states $\alpha\le\omega$ with equality
iff the items are tau-equivalent ($\lambda_m$ constant); both lie in $[0,1]$ since
each is a ratio of a nonnegative variance to a larger total variance. Under
tau-equivalence $\lambda_m\equiv\lambda$, $\sum_m\operatorname{Var}(x_m)=p(\lambda^2\sigma_\eta^2+\sigma_\varepsilon^2)$
and $\operatorname{Var}(X)=p^2\lambda^2\sigma_\eta^2+p\sigma_\varepsilon^2$; substituting and
simplifying gives $\alpha=\frac{p\lambda^2\sigma_\eta^2}{p\lambda^2\sigma_\eta^2+\sigma_\varepsilon^2}=\omega$.
\end{proof}
\emph{EI:} the value of $\hat\alpha$ (Edmondson reports $\alpha=0.82$, which by
Proposition~\ref{prop:alpha} certifies $\ge 0.82$ true-score variance) and the
test--retest correlation across two administrations (the population quantity is
$\operatorname{Corr}(\theta^{(\mathrm{PS})}_{j,t_1},\theta^{(\mathrm{PS})}_{j,t_2})$, which is
high iff $H(\mathrm{PS})$ is enduring, Section~\ref{sec:concept}).

\subsection{Aggregation to the team level}\label{sec:agg}
Define, for a one-way random-effects ANOVA of $s_{ij}$ on team $j$ (balanced
$n_j\equiv k$ for clean formulae), the between- and within-team mean squares
$\mathrm{MSB},\mathrm{MSW}$ and
\begin{equation}\label{eq:icc}
\mathrm{ICC}(1)=\frac{\mathrm{MSB}-\mathrm{MSW}}{\mathrm{MSB}+(k-1)\mathrm{MSW}},
\qquad
\mathrm{ICC}(2)=\frac{\mathrm{MSB}-\mathrm{MSW}}{\mathrm{MSB}} .
\end{equation}

\begin{proposition}[Spearman--Brown link]\label{prop:sb}
$\displaystyle \mathrm{ICC}(2)=\frac{k\,\mathrm{ICC}(1)}{1+(k-1)\,\mathrm{ICC}(1)}$,
and $\mathrm{ICC}(1)\to\dfrac{\sigma^2_{\theta,\mathrm{PS}}}
{\sigma^2_{\theta,\mathrm{PS}}+\sigma^2_{u,\mathrm{PS}}/p}$ in probability as
$J\to\infty$ (its population value is the fraction of total score variance that is
between-team).
\end{proposition}
\begin{proof}
Write $a=\mathrm{ICC}(1)$. From \eqref{eq:icc}, $a(\mathrm{MSB}+(k-1)\mathrm{MSW})=\mathrm{MSB}-\mathrm{MSW}$,
so $\mathrm{MSW}(1+(k-1)a)=\mathrm{MSB}(1-a)$, i.e.\
$\mathrm{MSW}/\mathrm{MSB}=(1-a)/(1+(k-1)a)$. Hence
$\mathrm{ICC}(2)=1-\mathrm{MSW}/\mathrm{MSB}=1-\frac{1-a}{1+(k-1)a}
=\frac{(1+(k-1)a)-(1-a)}{1+(k-1)a}=\frac{ka}{1+(k-1)a}$,
the Spearman--Brown formula for a mean of $k$ exchangeable raters. For the limit,
under \eqref{eq:cfa} the expected mean squares are
$\mathbb{E}[\mathrm{MSW}]=\sigma^2_w$ and
$\mathbb{E}[\mathrm{MSB}]=\sigma^2_w+k\sigma^2_b$ with
$\sigma^2_b=\sigma^2_{\theta,\mathrm{PS}}$ and
$\sigma^2_w=\sigma^2_{u,\mathrm{PS}}/p+\text{(error)}$; by the LLN
$\mathrm{MSW},\mathrm{MSB}$ converge to these expectations, and the continuous-mapping
theorem applied to \eqref{eq:icc} gives the stated limit. \end{proof}

This justifies aggregation: when $\mathrm{ICC}(1)$ is non-negligible and
$\mathrm{ICC}(2)$ (the reliability of the \emph{team mean}) is adequate, $\bar s_j$
reliably estimates $\theta^{(\mathrm{PS})}_j$, so the team mean is a meaningful
team-level variable.

\begin{proposition}[Within-group agreement]\label{prop:rwg}
For a single item with observed within-team variance $S_X^2$ on a scale whose
maximum-disagreement (uniform-null) variance is $\sigma^2_{EU}$, James--Demaree--Wolf
$r_{WG}=1-S_X^2/\sigma^2_{EU}$ satisfies $r_{WG}=1$ iff members agree exactly,
$r_{WG}=0$ iff within-team variance equals the no-agreement benchmark, and
$r_{WG}\in(0,1]$ iff $0\le S_X^2<\sigma^2_{EU}$. For $p$ essentially parallel items
the multi-item version
$r_{WG(p)}=\dfrac{p(1-\bar S^2/\sigma^2_{EU})}{p(1-\bar S^2/\sigma^2_{EU})+\bar S^2/\sigma^2_{EU}}$
is increasing in $p$.
\end{proposition}
\begin{proof}
The single-item claims are immediate from $r_{WG}=1-S_X^2/\sigma^2_{EU}$: it equals
$1$ iff $S_X^2=0$ (exact agreement), equals $0$ iff $S_X^2=\sigma^2_{EU}$, and lies
in $(0,1]$ iff $0\le S_X^2<\sigma^2_{EU}$. For the multi-item index put
$t=1-\bar S^2/\sigma^2_{EU}\in[0,1]$; then $r_{WG(p)}=pt/(pt+(1-t))$. Its derivative
in $p$ has numerator $t(pt+1-t)-pt\cdot t=t(1-t)\ge0$, so $r_{WG(p)}$ is
nondecreasing in $p$ (strictly increasing when $0<t<1$). \end{proof}
\emph{EI:} that the median $r_{WG}\ge 0.70$, $\mathrm{ICC}(1)$ and $\mathrm{ICC}(2)$
exceed conventional thresholds in the data (Edmondson reports values supporting
aggregation). Propositions~\ref{prop:sb}--\ref{prop:rwg} guarantee these indices
\emph{mean} what aggregation requires.

\subsection{(2b) Convergent validity}
Convergent validity is the requirement $\phi_{\mathrm{PS},c'}$ \emph{large and
positive} for theoretically related $c'$ (PS should correlate substantially with
trust, cohesion, satisfaction, efficacy). Within \eqref{eq:cfa} this is the claim
that the estimated factor correlations $\hat\phi_{\mathrm{PS},c'}$ are bounded away
from $0$. The mathematical content is the disattenuation identity below; the
magnitudes are EI.

\begin{lemma}[Disattenuation]\label{lem:disatt}
If $\hat\rho$ is the observed correlation between the PS and $c'$ \emph{scale
scores} with reliabilities $\rho_{\mathrm{PS}},\rho_{c'}$, then the true
factor correlation is $\phi_{\mathrm{PS},c'}=\hat\rho/\sqrt{\rho_{\mathrm{PS}}\rho_{c'}}$
in the population.
\end{lemma}
\begin{proof}
By \eqref{eq:cfa} each scale score $S_c=T_c+E_c$ with true part $T_c$ and
independent error $E_c$, $\operatorname{Var}(T_c)/\operatorname{Var}(S_c)=\rho_c$. Then
$\operatorname{Cov}(S_{\mathrm{PS}},S_{c'})=\operatorname{Cov}(T_{\mathrm{PS}},T_{c'})$, so
$\operatorname{Corr}(S_{\mathrm{PS}},S_{c'})=\frac{\operatorname{Cov}(T_{\mathrm{PS}},T_{c'})}{\sqrt{\operatorname{Var}(S_{\mathrm{PS}})\operatorname{Var}(S_{c'})}}
=\phi_{\mathrm{PS},c'}\sqrt{\rho_{\mathrm{PS}}\rho_{c'}}$. Solving gives the claim. \end{proof}

\subsection{(2c) Discriminant validity via nested CFA}\label{sec:disc}
This is the heart of Task~2. The discriminant-validity claim for the pair
$(\mathrm{PS},c')$ is $|\phi_{\mathrm{PS},c'}|<1$, established by showing a
\emph{two-factor} model fits significantly better than the \emph{merged
one-factor} model.

\begin{definition}[The two competing models]
Let $\mathcal M_2(c')$ be the CFA model in which the PS items load on $\eta^{(\mathrm{PS})}$,
the $c'$ items on $\eta^{(c')}$, with free correlation $\phi_{\mathrm{PS},c'}\in(-1,1)$.
Let $\mathcal M_1(c')$ be the same model with the single equality constraint
$\phi_{\mathrm{PS},c'}=1$ (a single common factor for the merged item set).
\end{definition}

\begin{lemma}[Nesting]\label{lem:nest}
$\mathcal M_1(c')$ is nested in $\mathcal M_2(c')$: its parameter space is the
intersection of that of $\mathcal M_2(c')$ with the hyperplane
$\{\phi_{\mathrm{PS},c'}=1\}$, and the difference in numbers of free parameters is
exactly $1$.
\end{lemma}
\begin{proof}
Both models share identical loading, intercept, and unique-variance parameters and
the same observed-variable set. The only parameter free in $\mathcal M_2$ but fixed
in $\mathcal M_1$ is $\phi_{\mathrm{PS},c'}$. Setting it to its boundary-interior
value $1$ produces a model whose implied covariance matrix
$\Sigma(\vartheta)$ is the restriction of $\mathcal M_2$'s to that hyperplane.
Hence $\mathcal M_1\subset\mathcal M_2$ as parameter sets and
$\dim\mathcal M_2-\dim\mathcal M_1=1$, i.e.\ $\mathrm{df}_1-\mathrm{df}_2=1$. \end{proof}

\begin{lemma}[Identifiability]\label{lem:ident}
If each latent factor has $p\ge 3$ indicators with nonzero loadings, errors
uncorrelated, and one loading per factor fixed (or factor variance fixed to $1$)
for scale-setting, then $\mathcal M_2(c')$ is (locally) identified; in particular
$\phi_{\mathrm{PS},c'}$ is a function of the population covariance matrix.
\end{lemma}
\begin{proof}
This is the standard three-indicator rule. With factor variances normalised to
$1$, for a single factor with indicators $1,2,3$ the off-diagonal covariances give
$\sigma_{12}=\lambda_1\lambda_2$, $\sigma_{13}=\lambda_1\lambda_3$,
$\sigma_{23}=\lambda_2\lambda_3$, whence
$\lambda_1^2=\sigma_{12}\sigma_{13}/\sigma_{23}$ and analogously $\lambda_2,\lambda_3$
are recovered (sign fixed by a marker), then unique variances from the diagonal.
With both factors so identified, any cross-loading covariance
$\sigma_{x y}=\lambda_x\lambda_y\,\phi_{\mathrm{PS},c'}$ ($x$ a PS indicator, $y$ a
$c'$ indicator) yields $\phi_{\mathrm{PS},c'}=\sigma_{xy}/(\lambda_x\lambda_y)$, a
well-defined function of $\Sigma$. Local identification then follows from the
column rank of the Jacobian $\partial\mathrm{vech}\,\Sigma/\partial\vartheta$ being
full at the true parameter. \end{proof}

\begin{theorem}[Validity of the discriminant-validity test]\label{thm:diff}
Let $\hat\Sigma$ be the sample covariance matrix from $N$ independent teams and let
$F_{\mathrm{ML}}(\vartheta)=\log|\Sigma(\vartheta)|+\operatorname{tr}(\hat\Sigma\Sigma(\vartheta)^{-1})
-\log|\hat\Sigma|-q$ be the normal-theory ML discrepancy ($q$ the number of observed
variables). Let $\hat\vartheta_1,\hat\vartheta_2$ minimise $F_{\mathrm{ML}}$ over
$\mathcal M_1(c'),\mathcal M_2(c')$ and set
$T=N\big(F_{\mathrm{ML}}(\hat\vartheta_1)-F_{\mathrm{ML}}(\hat\vartheta_2)\big)
=\chi^2_{\mathcal M_1}-\chi^2_{\mathcal M_2}$.
Assume $\mathcal M_2(c')$ is correctly specified and identified
(Lemma~\ref{lem:ident}) and the data are i.i.d.\ with finite fourth moments. Then:
\begin{enumerate}
\item[(i)] If the constraint $\phi_{\mathrm{PS},c'}=1$ holds in the population
(merged constructs), $T\xrightarrow{d}\chi^2_1$.
\item[(ii)] If $\phi_{\mathrm{PS},c'}<1$ in the population (distinct constructs),
$T\to\infty$ in probability; equivalently the noncentrality of $T$ grows linearly
in $N$, so the test rejecting $\mathcal M_1$ when $T>\chi^2_{1,1-\alpha}$ has power
$\to1$.
\end{enumerate}
Consequently ``$\mathcal M_2$ fits significantly better than $\mathcal M_1$''
($T>\chi^2_{1,0.95}=3.84$) is, asymptotically, a level-$\alpha$ test of the
discriminant-validity hypothesis $H_1:\phi_{\mathrm{PS},c'}<1$ against the
merge hypothesis $H_0:\phi_{\mathrm{PS},c'}=1$.
\end{theorem}
\begin{proof}
By Lemma~\ref{lem:nest} the two models differ by the single smooth equality
constraint $g(\vartheta)=\phi_{\mathrm{PS},c'}-1=0$, and by Lemma~\ref{lem:ident}
$\mathcal M_2$ is identified, so $\hat\vartheta_2$ is consistent and asymptotically
normal with the usual ML covariance $\frac1N\,\mathcal I(\vartheta_0)^{-1}$, where
$\mathcal I$ is the (nonsingular) Fisher information of the normal-theory model.
The statistic $T$ is exactly the likelihood-ratio statistic for testing
$g(\vartheta)=0$ within $\mathcal M_2$ (the SEM fit function $NF_{\mathrm{ML}}$ is,
up to an additive constant independent of $\vartheta$, $-2$ times the multivariate
normal log-likelihood). 

(i) Under $H_0$ the constraint is a single regular restriction of an identified
model satisfying the LR regularity conditions (here $\phi=1$ is reached as an
interior limit along the admissible region; the score for $g$ is non-degenerate by
Lemma~\ref{lem:ident}); the classical Wilks theorem gives $T\xrightarrow{d}\chi^2_r$
with $r=\dim\mathcal M_2-\dim\mathcal M_1=1$ by Lemma~\ref{lem:nest}.

(ii) Under a fixed alternative $\phi_{\mathrm{PS},c'}=\phi_0<1$, $\mathcal M_1$ is
misspecified: its best-fitting population covariance $\Sigma(\vartheta_1^\ast)$
differs from the true $\Sigma_0$, so $F_{\mathrm{ML}}(\vartheta_1^\ast)>0
=F_{\mathrm{ML}}(\vartheta_2^\ast)$ (the discrepancy is zero iff the model
covariance equals the truth, and $\Sigma(\vartheta_1^\ast)\neq\Sigma_0$ because no
member of $\mathcal M_1$ reproduces a between-factor correlation $<1$ while
matching the within-factor structure---formally, by Lemma~\ref{lem:ident},
$\phi_{\mathrm{PS},c'}$ is a continuous function of $\Sigma$ equal to $\phi_0<1$ at
$\Sigma_0$, but identically $1$ on $\mathcal M_1$). By consistency of the ML
discrepancy, $F_{\mathrm{ML}}(\hat\vartheta_1)\xrightarrow{p}F_{\mathrm{ML}}(\vartheta_1^\ast)=:\delta>0$
and $F_{\mathrm{ML}}(\hat\vartheta_2)\xrightarrow{p}0$, so
$T=N(F_{\mathrm{ML}}(\hat\vartheta_1)-F_{\mathrm{ML}}(\hat\vartheta_2))$ behaves like
$N\delta\to\infty$; equivalently the local-alternative noncentrality parameter is
$\lambda=N\delta$, giving power $\to1$. The displayed level/threshold statements
follow from (i)--(ii). \end{proof}

\begin{remark}[Confidence-interval form]
An equivalent and often-reported criterion is that the Wald confidence interval
$\hat\phi_{\mathrm{PS},c'}\pm z_{1-\alpha/2}\,\widehat{\mathrm{SE}}(\hat\phi_{\mathrm{PS},c'})$
excludes $1$; by the asymptotic normality of $\hat\phi_{\mathrm{PS},c'}$ established
in the proof of Theorem~\ref{thm:diff}, this Wald test is asymptotically equivalent
to the LR test for $H_0:\phi=1$ (Wald--LR--score equivalence under the regularity
conditions). The Fornell--Larcker criterion
$\phi_{\mathrm{PS},c'}^2<\min(\mathrm{AVE}_{\mathrm{PS}},\mathrm{AVE}_{c'})$, where
$\mathrm{AVE}_c=\frac1p\sum_m(\lambda^{(c)}_m)^2/[(\lambda^{(c)}_m)^2+\operatorname{Var}(\varepsilon^{(c)}_m)]$
is the average variance extracted, is a strictly stronger sufficient condition for
$|\phi_{\mathrm{PS},c'}|<1$, since $\mathrm{AVE}_c\le1$ forces
$\phi^2<\min\mathrm{AVE}\le1$.
\end{remark}

\emph{EI:} the actual $\chi^2$-difference values for each $c'\in\{\mathrm{COH,TR,SAT,EFF}\}$.
In Edmondson's data each two-factor model fits significantly better than the merged
model (each $\chi^2$ difference exceeds $3.84$ on $1\,$df), i.e.\ $H_0$ is rejected
for every related construct. Given these EI, Theorem~\ref{thm:diff} delivers
$\phi_{\mathrm{PS},c'}<1$ for all four pairs.

\subsection{(2d) Nomological / incremental validity}\label{sec:nomo}
Let $Y_j$ be a team outcome (e.g.\ learning behaviour, error detection,
performance). Consider the hierarchical population regressions
\begin{align}
\text{(R0)}\quad & Y_j=\beta_0+\textstyle\sum_{c'\neq\mathrm{PS}}\gamma_{c'}\theta^{(c')}_j+e_j,\\
\text{(R1)}\quad & Y_j=\beta_0'+\beta_{\mathrm{PS}}\theta^{(\mathrm{PS})}_j
+\textstyle\sum_{c'\neq\mathrm{PS}}\gamma'_{c'}\theta^{(c')}_j+e_j'.
\end{align}

\begin{proposition}[Incremental validity criterion]\label{prop:incr}
$\theta^{(\mathrm{PS})}_j$ has predictive validity for $Y$ \emph{above and beyond}
the related constructs iff the partial coefficient $\beta_{\mathrm{PS}}\neq0$,
equivalently iff the population increment in explained variance
$\Delta R^2=R^2_{\mathrm{R1}}-R^2_{\mathrm{R0}}>0$. Moreover, if
$|\phi_{\mathrm{PS},c'}|<1$ for all $c'$ (discriminant validity holds) then PS is
not collinear with the block $\{\theta^{(c')}\}$, so $\beta_{\mathrm{PS}}$ is
identified and the increment $\Delta R^2$ is testable by a partial $F$-test.
\end{proposition}
\begin{proof}
By the Frisch--Waugh--Lovell theorem, $\beta_{\mathrm{PS}}$ equals the simple
regression coefficient of $Y$ on the residual of $\theta^{(\mathrm{PS})}$ after
projecting out the block $Z=\{\theta^{(c')}\}_{c'\neq\mathrm{PS}}$; call this
residual $\theta^{(\mathrm{PS})}_{\perp}$. Then
$\Delta R^2=\operatorname{Corr}(Y,\theta^{(\mathrm{PS})}_\perp)^2\,(1-R^2_{Y\sim Z})$ up to the
usual normalisation, which is $>0$ iff
$\operatorname{Cov}(Y,\theta^{(\mathrm{PS})}_\perp)\neq0$ iff $\beta_{\mathrm{PS}}\neq0$. The
residual $\theta^{(\mathrm{PS})}_\perp$ has positive variance iff
$\theta^{(\mathrm{PS})}$ is not a linear combination of $Z$; since the Gram matrix
of $\{\theta^{(\mathrm{PS})}\}\cup Z$ is $\Phi$ and discriminant validity gives every
off-diagonal $|\phi_{\mathrm{PS},c'}|<1$, $\Phi$ is nonsingular under the standard
assumption that the related block is itself non-collinear, so
$\operatorname{Var}(\theta^{(\mathrm{PS})}_\perp)>0$ and $\beta_{\mathrm{PS}}$ is identified. The
sample increment is tested by the partial $F$-statistic
$F=\frac{(\mathrm{RSS}_0-\mathrm{RSS}_1)/1}{\mathrm{RSS}_1/(N-\dim-1)}$, which under
$H_0:\beta_{\mathrm{PS}}=0$ has an $F_{1,N-\dim-1}$ distribution and diverges under
$\beta_{\mathrm{PS}}\neq0$. \end{proof}
\emph{EI:} that $\hat\beta_{\mathrm{PS}}>0$ and the increment is significant for
the outcomes (Edmondson finds PS predicts team learning behaviour, and learning
behaviour mediates PS's effect on performance, controlling for the related
constructs). Given these EI, Proposition~\ref{prop:incr} certifies nomological
validity.

\section{Task 3: Verdict and required further work}\label{sec:verdict}

\begin{theorem}[Construct-validity verdict]\label{thm:verdict}
Suppose the following empirical inputs hold for the population of teams (each is the
observed pattern in Edmondson 1999 and is the kind of statement only data can
establish):
\begin{itemize}
\item[(EI1)] $\hat\alpha\ge0.70$ and adequate test--retest stability for PS;
\item[(EI2)] $r_{WG}$, $\mathrm{ICC}(1)$, $\mathrm{ICC}(2)$ above conventional
aggregation thresholds;
\item[(EI3)] for every $c'\in\{\mathrm{COH,TR,SAT,EFF}\}$ the $\chi^2$-difference
$\chi^2_{\mathcal M_1(c')}-\chi^2_{\mathcal M_2(c')}>3.84$;
\item[(EI4)] $\hat\beta_{\mathrm{PS}}\neq0$ with significant $\Delta R^2$ for at
least one substantive outcome, controlling for the related constructs.
\end{itemize}
Then team psychological safety is \emph{both conceptually and empirically distinct}
(Definition~\ref{def:distinct}) from group cohesiveness, interpersonal trust,
work-group satisfaction, and team efficacy.
\end{theorem}
\begin{proof}
Conceptual distinctness is Proposition~\ref{prop:concept} (unconditional). For
empirical distinctness we must show $|\phi_{\mathrm{PS},c'}|<1$ for each $c'$.
By (EI2) and Propositions~\ref{prop:sb}--\ref{prop:rwg} aggregation is justified, so
$\theta^{(\mathrm{PS})}_j$ is a well-defined team-level variable measured by
$\bar s_j$ with reliability $\mathrm{ICC}(2)$; by (EI1) and
Proposition~\ref{prop:alpha} the PS scale carries true-score variance, so the
factor $\eta^{(\mathrm{PS})}$ in \eqref{eq:cfa} is non-degenerate and the CFA models
$\mathcal M_1,\mathcal M_2$ are identified (Lemma~\ref{lem:ident}). By (EI3) and
Theorem~\ref{thm:diff}(ii)--threshold, for each $c'$ the test rejects
$H_0:\phi_{\mathrm{PS},c'}=1$ at level $0.05$, i.e.\ the data are consistent only
with $\phi_{\mathrm{PS},c'}<1$; hence PS is empirically distinct from each of the
four. (EI4) with Proposition~\ref{prop:incr} additionally certifies that PS adds
predictive information beyond the four constructs, ruling out the
``redundant-predictor'' interpretation. Combining, PS satisfies both clauses of
Definition~\ref{def:distinct} relative to every listed construct. \end{proof}

\begin{remark}[Honest reading of the verdict]
Theorem~\ref{thm:verdict} is a conditional theorem: it shows that the standard
construct-validation criteria, \emph{when the data exhibit patterns
(EI1)--(EI4)}, deductively entail distinctness. The antecedents are empirical and
were verified in Edmondson's field study; they are not provable from mathematics
alone, and a different population could in principle violate them. What is proved
unconditionally here is (a) the conceptual distinctness
(Proposition~\ref{prop:concept}) and (b) the \emph{validity of every inferential
step}---the aggregation indices mean what they must (Props.\
\ref{prop:sb}--\ref{prop:rwg}), the reliability bound holds
(Prop.~\ref{prop:alpha}), the merged-vs.\ two-factor comparison is a correctly
sized, consistent test of $\phi=1$ (Thm.~\ref{thm:diff}), and the
incremental-validity criterion is identified and testable (Prop.~\ref{prop:incr}).
\end{remark}

\subsection*{Required further conceptual and empirical work}
The construct's meaning, boundary conditions, antecedents, and consequences still
require:
\begin{enumerate}
\item \textbf{Generalisability (external validity).} (EI1)--(EI4) were established
in particular samples; replication across industries, cultures, and team types is
needed, since $\Phi$ and the $\lambda$'s are population-specific. \emph{Measurement
invariance} (equality of $\nu^{(c)}_m,\lambda^{(c)}_m$ across groups) must be tested
before cross-group comparisons of $\theta^{(\mathrm{PS})}$ are meaningful; this is
itself a nested-model comparison of the type in Theorem~\ref{thm:diff}.
\item \textbf{Dyadic-to-collective bridge.} Because $R(\mathrm{TR})$ is dyadic
(Section~\ref{sec:concept}), a formal model linking the dyadic trust array to the
shared belief $\theta^{(\mathrm{PS})}$ (e.g.\ via a social-network aggregation
operator) would sharpen the PS/TR boundary beyond the scalar $\phi_{\mathrm{PS,TR}}$.
\item \textbf{Causal nomological net.} Theorem~\ref{thm:verdict}(EI4) is
associational. Establishing antecedents (leader inclusiveness, team structure) and
consequences (learning, performance) as \emph{causal} requires longitudinal or
(quasi-)experimental designs and mediation models with identification assumptions
beyond the cross-sectional CFA.
\item \textbf{Boundary conditions / moderators.} The sign and size of PS's effect
may depend on task interdependence and accountability; mapping these moderators
refines the construct's consequences.
\item \textbf{Higher reliability of $\mathrm{ICC}(2)$.} By
Proposition~\ref{prop:sb}, $\mathrm{ICC}(2)$ rises with team size $k$; small teams
need either more items or model-based latent aggregation to keep team-mean
reliability high.
\end{enumerate}

\section*{Conclusion}
Conceptually, team psychological safety is distinct from cohesiveness, trust,
satisfaction, and efficacy on at least one of level, referent, or time horizon
(Proposition~\ref{prop:concept}). Empirically, under the validation patterns that
hold in field data, the formal criteria entail that it is a distinct latent
construct: aggregation is justified (Props.~\ref{prop:sb}--\ref{prop:rwg}), the
scale is reliable (Prop.~\ref{prop:alpha}), each merged one-factor model fits
significantly worse than the corresponding two-factor model
(Theorem~\ref{thm:diff}), and PS predicts outcomes incrementally
(Prop.~\ref{prop:incr}). Hence team psychological safety is an empirically distinct
construct (Theorem~\ref{thm:verdict}), with the further work above needed to
delimit its generality, causal role, and boundary conditions.
\end{document}
